\section{Methodology}

In this section, we present the overall architecture of the proposed Unified Dual-Mask Physical Model for non-Lambertian light field depth estimation. Different from previous works that rely on implicit feature learning [2] or continuous angular spectrum analysis [2], our framework explicitly integrates physical reflectance priors into a deep neural network. The architecture is designed to decouple the complex light field signal into distinct physical components and process them through specialized pathways. As illustrated in Fig. 1, the model comprises a physical decomposition frontend, a dual-branch feature extraction backbone, and a mask-guided fusion backend. This topology enables the network to adaptively route and process Lambertian and non-Lambertian regions, yielding accurate depth maps and domain classification masks simultaneously.

The network takes a $9 \times 9$ angular light field image patch, denoted as $I(u,v)$, where $u$ and $v$ represent the angular coordinates. To address the limitations of frequency-domain methods that require dense angular sampling [2], we introduce the Three-Layer Angular Signal Decomposition module. The motivation is to explicitly disentangle the raw light field signal into physically meaningful components without relying on continuous angular resolution. Specifically, the module decomposes the angular signal into three independent layers: a direct current (DC) component representing ambient illumination, a linear component parameterized by coefficients $a$ and $b$ representing disparity (i.e., epipolar plane image (EPI) slopes), and a residual component $\varepsilon(u,v)$ capturing material properties and reflectance types. The physical decomposition is formulated as:
\begin{equation}
\label{eq:eq2}
I(u,v) = DC + a \cdot u + b \cdot v + \varepsilon(u,v)
\end{equation}
where $DC$, $a$, and $b$ are estimated via least-squares fitting across the angular dimensions. Concurrently, the module computes highly discriminative geometric features, such as the angular gradient $\nabla_{angular} I$, directly from the residual term. This explicit physical decoupling ensures that the subsequent branches receive purified signals, substantially improving the robustness of feature extraction in discrete $9 \times 9$ angular sampling scenarios.

Following the physical decomposition, the data flows into a Dual-Branch Architecture designed to handle the distinct geometric characteristics of different surface types. The motivation for this parallel design is that a single unified branch struggles to simultaneously model the strict photo-consistency of Lambertian surfaces and the complex view-dependent variations of non-Lambertian surfaces [2]. The upper pathway, the Lambertian EPI Branch, receives the linear coefficients $(a,b)$ and EPI structural features. It leverages the linear EPI line structures and incorporates a piecewise smooth continuous depth prior via a Planar Regularized Sampling Operator (PRSO) to refine depth estimation in diffuse regions. Conversely, the lower pathway, the Non-Lambertian Geometric Branch, processes the residual component $\varepsilon(u,v)$ and the extracted geometric features. Instead of relying on disrupted EPI line slopes, this branch detects X-shape patterns and bimodal distributions within the EPIs [2]. It employs geometric convolutions to extract depth and material cues from non-Lambertian regions, augmented by Random Tone-mapping Augmentation (RTA) to mitigate data scarcity. By separating the processing logic, the dual-branch design effectively prevents the interference between diffuse and specular feature representations.

To integrate the complementary features from the dual branches, the architecture employs a Dual-Mask Routing \& Fusion module. The primary motivation is to avoid suboptimal scene-level binary cross-entropy classification and instead achieve precise pixel-level adaptive fusion. The module takes the feature maps from both branches and the highly discriminative angular gradient (which achieves a Cohen's d of 0.983 between domains) as inputs. First, the angular gradient is fed into a multi-layer perceptron followed by a sigmoid activation to generate a pixel-level domain mask $M$:
\begin{equation}
\label{eq:eq3}
M = \sigma(\text{MLP}(\nabla_{angular} I))
\end{equation}
Simultaneously, the output features from the Lambertian and non-Lambertian branches are projected to a unified channel dimension using a FUSE\_DIM projection layer (with the channel size set to 96) to balance the feature representations. The final depth map $D$ is then computed through a mask-weighted fusion operation:
\begin{equation}
\label{eq:eq4}
D = M \odot D_L + (1 - M) \odot D_{NL}
\end{equation}
where $D_L$ and $D_{NL}$ denote the projected depth features from the Lambertian and non-Lambertian branches, respectively, and $\odot$ represents the element-wise multiplication. Ultimately, the network outputs the high-precision Depth Map and the Domain Mask. This mask-guided fusion mechanism ensures that each pixel's depth is derived from the most physically appropriate branch, significantly enhancing the overall estimation accuracy in complex mixed scenes.

\subsection{Three-Layer Angular Signal Decomposition}

The primary objective of the Three-Layer Angular Signal Decomposition module is to explicitly disentangle the raw light field signal into physically meaningful components without relying on continuous angular resolution. Conventional light field depth estimation methods predominantly rely on Epipolar Plane Image (EPI) slopes [2] or continuous angular spectrum analysis [2], which inherently assume Lambertian reflectance. When encountering non-Lambertian surfaces, such as specular highlights or complex reflections, the angular signal undergoes severe non-linear distortions, causing these conventional methods to fail. Different from previous works that implicitly learn feature representations [2] or require dense angular sampling for frequency-domain analysis [2], our module explicitly formulates the angular signal decomposition as a physical parameter estimation problem. This design isolates the Lambertian components for accurate depth extraction while preserving the non-Lambertian residuals for material and geometric structure analysis.

Let the input to this module be a $9 \times 9$ angular light field image patch, denoted as $I(u,v)$, where $u$ and $v$ represent the discrete angular coordinates, $u, v \in \{-4, -3, \dots, 4\}$. The physical reflectance model parameterizes the angular intensity variation as a superposition of ambient illumination, parallax-induced linear shifts, and material-specific non-linear residuals. Mathematically, the angular signal is modeled as:
\begin{equation}
\label{eq:eq5}
I(u,v) = I_{dc} + a \cdot u + b \cdot v + \varepsilon(u,v)
\end{equation}
where $I_{dc}$ represents the direct current (DC) component corresponding to the ambient illumination or albedo, $a$ and $b$ are the linear coefficients along the $u$ and $v$ angular axes encoding the parallax (and thus depth) information, and $\varepsilon(u,v)$ denotes the residual component capturing non-Lambertian effects such as specularities and complex reflections.

To solve for the physical parameters, we employ an orthogonal projection approach. Given the symmetric distribution of the angular coordinates around zero, the basis functions $\phi_0(u,v) = 1$, $\phi_1(u,v) = u$, and $\phi_2(u,v) = v$ are mutually orthogonal. This orthogonality allows for the closed-form analytical derivation of the parameters without iterative optimization, ensuring computational efficiency. The DC component $I_{dc}$ is computed as the mean intensity across the angular domain:
\begin{equation}
\label{eq:eq6}
I_{dc} = \frac{1}{N} \sum_{u=-4}^{4} \sum_{v=-4}^{4} I(u,v)
\end{equation}
where $N = 81$ is the total number of angular samples. The linear coefficients $a$ and $b$, which dictate the EPI slopes and correspond to the horizontal and vertical disparities, are derived by projecting the signal onto the respective angular axes:
\begin{equation}
\label{eq:eq7}
a = \frac{\sum_{u=-4}^{4} \sum_{v=-4}^{4} u \cdot I(u,v)}{\sum_{u=-4}^{4} \sum_{v=-4}^{4} u^2}, \quad b = \frac{\sum_{u=-4}^{4} \sum_{v=-4}^{4} v \cdot I(u,v)}{\sum_{u=-4}^{4} \sum_{v=-4}^{4} v^2}
\end{equation}

Subsequently, the residual component $\varepsilon(u,v)$ is obtained by subtracting the estimated DC and linear components from the original angular signal:
\begin{equation}
\label{eq:eq8}
\varepsilon(u,v) = I(u,v) - \left( I_{dc} + a \cdot u + b \cdot v \right)
\end{equation}
This residual map inherently encapsulates the high-frequency non-Lambertian variations. To further extract discriminative geometric cues, we compute the angular gradient $\nabla_{angular} I$ directly from $\varepsilon(u,v)$ using Sobel operators along the $u$ and $v$ dimensions. The resulting feature tensors, including $I_{dc}$, the disparity coefficients $(a, b)$, the residual map $\varepsilon$, and the angular gradient $\nabla_{angular} I$, are then concatenated and forwarded to the subsequent Dual-Branch Architecture. This explicit decomposition not only provides physically grounded initial features but also significantly reduces the learning burden on the downstream deep neural networks by eliminating the need to implicitly disentangle mixed reflectance properties.